\begin{lstlisting}[
    caption={Trénovanie modelu YOLO11},
  label={lst:python},
  language=python,
  style=code-listing
]
from ultralytics import YOLO
import pandas as pd
import os

# =====================================================
# DATASET PATHS
# =====================================================

DEFAULT_DATASET = "/path to YAML"

BINARY_DATASET = "/path to bin YAML"

# =====================================================
# TRENING
# =====================================================

def run_training(
    model_name,
    dataset_path,
    run_name,
    augmentation=False
):

    print("\n====================================")
    print(f"Starting: {run_name}")
    print("====================================\n")

    # load model
    model = YOLO(model_name)

    # =============================================
    # TRAIN
    # =============================================

    metrics = model.train(

        # dataset
        data=dataset_path,

        # max epochs
        epochs=200,

        # EARLY STOPPING
        patience=40,

        # image size
        imgsz=640,

        # batch size
        batch=8,

        # =========================================
        # AUGMENTACIA
        # =========================================

        hsv_h=0.015 if augmentation else 0.0,
        hsv_s=0.7 if augmentation else 0.0,
        hsv_v=0.4 if augmentation else 0.0,

        degrees=10 if augmentation else 0.0,

        translate=0.1 if augmentation else 0.0,

        scale=0.5 if augmentation else 0.0,

        shear=10 if augmentation else 0.0,

        fliplr=0.0,
        flipud=0.0,

        # =========================================
        # SAVE
        # =========================================

        project="custom output path",

        name=run_name
    )



# =====================================================
# DEF DATASET
# =====================================================

run_training(

    model_name="yolo11s.pt",

    dataset_path=DEFAULT_DATASET,

    run_name="yolo11s_default_earlystop_p40",

    augmentation=False
)

# =====================================================
# BIN DATASET
# =====================================================

run_training(

    model_name="yolo11s.pt",

    dataset_path=BINARY_DATASET,

    run_name="yolo11s_binary_earlystop_p40",

    augmentation=False
)

# =====================================================
# AUG DATASET
# =====================================================

run_training(

    model_name="yolo11s.pt",

    dataset_path=DEFAULT_DATASET,

    run_name="yolo11s_augmented_earlystop_p40",

    augmentation=True
)

\end{lstlisting}

\begin{lstlisting}[
    caption={Kód použitý na binarizáciu datasetu},
  label={lst:python},
  language=python,
  style=code-listing
]

import cv2
import os
import shutil
from glob import glob

# =====================================================
# INPUT / OUTPUT PATHS
# =====================================================

# ORIGINAL DATASET
input_images = "/input_folder_imgs"
input_labels = "/input_folder_labels"

# NEW BINARY DATASET
output_images = "/output_folder_imgs"
output_labels = "/output_folder_labels"

# create folders
os.makedirs(output_images, exist_ok=True)
os.makedirs(output_labels, exist_ok=True)

# =====================================================
# PROCESS IMAGES
# =====================================================

for img_path in glob(input_images + "/*"):

    # =============================================
    # LOAD IMAGE
    # =============================================

    img = cv2.imread(img_path)

    # grayscale
    gray = cv2.cvtColor(img, cv2.COLOR_BGR2GRAY)

    # blur
    blur = cv2.GaussianBlur(gray, (5,5), 0)

    # =============================================
    # ADAPTIVE BINARIZATION
    # =============================================

    binary = cv2.adaptiveThreshold(

        blur,

        255,

        cv2.ADAPTIVE_THRESH_GAUSSIAN_C,

        cv2.THRESH_BINARY,

        21,

        15
    )

    # =============================================
    # REMOVE SMALL NOISE
    # =============================================

    kernel = cv2.getStructuringElement(
        cv2.MORPH_RECT,
        (2,2)
    )

    cleaned = cv2.morphologyEx(
        binary,
        cv2.MORPH_OPEN,
        kernel
    )

    cleaned = cv2.medianBlur(cleaned, 3)

    # =============================================
    # SAVE IMAGE
    # =============================================

    filename = os.path.basename(img_path)

    output_image_path = os.path.join(
        output_images,
        filename
    )

    cv2.imwrite(output_image_path, cleaned)

    # =============================================
    # COPY LABEL
    # =============================================

    label_name = os.path.splitext(filename)[0] + ".txt"

    input_label_path = os.path.join(
        input_labels,
        label_name
    )

    output_label_path = os.path.join(
        output_labels,
        label_name
    )

    # copy label if exists
    if os.path.exists(input_label_path):

        shutil.copy(
            input_label_path,
            output_label_path
        )

print("Binarization + label copy completed.")

\end{lstlisting}

\begin{lstlisting}[
    caption={Kód na spočítanie labelov pre štatistiku},
  label={lst:python},
  language=python,
  style=code-listing
]
import os
import json
from collections import defaultdict

def print_all_labels(json_dir):
# Dictionary to hold label counts
    label_counts = defaultdict(int)
    # Iterate over all JSON files in the directory
    for json_file in os.listdir(json_dir):
        if json_file.endswith(".json"):
            json_path = os.path.join(json_dir, json_file)
            with open(json_path, "r") as f:
                data = json.load(f)
            # Extract labels from the JSON file
            for shape in data.get("shapes", []):
                label = shape.get("label", "no-label")
                label_counts[label] += 1
    return label_counts

# Directory containing the JSON files
json_dir = "dataset/img:json"
# Find labels and their counts
labels = print_all_labels(json_dir)
# Print results
print("Labels and their counts:")
for label, count in labels.items():
    print(f"{label}: {count}")

\end{lstlisting}

\begin{lstlisting}[
  caption={Trénovanie YOLO11},
  label={lst:python},
  language=python,
  style=code-listing
]
import gc
import torch
from ultralytics import YOLO

# Context manager for YOLO model to ensure proper cleanup
class YOLOModel:
    def __init__(self, model_path):
        self.model_path = model_path
        self.model = None

    def __enter__(self):
        self.model = YOLO(self.model_path)
        return self.model

    def __exit__(self, exc_type, exc_val, exc_tb):
        del self.model
        gc.collect()
        if torch.cuda.is_available():
            torch.cuda.empty_cache()
        return False  # Don't suppress exceptions

# --- Update these paths as needed ---
MODEL_PATH = "yolo11x.pt"  # alebo "yolo11x.pt", "yolo11n.pt", etc.
DATASET_YAML = "dataset/YOLODataset/dataset.yaml"

def main():
    # Train the model using the context manager
    with YOLOModel(MODEL_PATH) as model:
        results = model.train(
            data=DATASET_YAML,
            epochs=100,
            imgsz=640,
            batch=8,
            cache=False,
            workers=0,
            device=0
        )
    # Optional cleanup (redundant, but safe)
    gc.collect()
    if torch.cuda.is_available():
        torch.cuda.empty_cache()

if __name__ == "__main__":
    main()

\end{lstlisting}
\begin{lstlisting}[
    caption={Detekcia vo viacerých obrázkoch},
  label={lst:python},
  language=python,
  style=code-listing
]
import os
import cv2
import numpy as np
from ultralytics import YOLO

# Paths
MODEL_PATH = "RegularNEW100epchSmall/weights/best.pt" #cesta k váham
INPUT_FOLDER = "dataset/YOLODataset/images/val" #cesta k input folderu s obrazkami
OUTPUT_FOLDER = "val X" #cesta k output folderu

# Supported image extensions
VALID_EXTENSIONS = (".jpg", ".jpeg", ".png")

# Create the output folder if it doesn't exist
os.makedirs(OUTPUT_FOLDER, exist_ok=True)

# Load the YOLO model
model = YOLO(MODEL_PATH)

def is_image_file(filename):
    return filename.lower().endswith(VALID_EXTENSIONS)

def process_image(input_path, output_path):
    # Read the image
    image = cv2.imread(input_path)
    if image is None:
        print(f"Warning: Could not read {input_path}")
        return

    # Run YOLO prediction
    results = model.predict(image, save=False)

    # Annotate the image
    annotated_pil = results[0].plot(font_size=60, pil=True)
    annotated_bgr = cv2.cvtColor(np.array(annotated_pil), cv2.COLOR_RGB2BGR)

    # Save the annotated image
    cv2.imwrite(output_path, annotated_bgr)
    print(f"Processed: {os.path.basename(input_path)}")

def main():
    for filename in os.listdir(INPUT_FOLDER):
        if is_image_file(filename):
            input_path = os.path.join(INPUT_FOLDER, filename)
            output_path = os.path.join(OUTPUT_FOLDER, filename)
            process_image(input_path, output_path)

if __name__ == "__main__":
    main()
\end{lstlisting}

\begin{lstlisting}[
  caption={Kód použitý pre binarizáciu a augmentáciu datasetu},
  label={lst:python},
  language=python,
  style=code-listing
]

import os
import cv2
import albumentations as A
import shutil

# === Setup Paths ===
ORIGINAL_DATASET = "dataset/YOLODataset"
BINARIZED_DATASET = "dataset/YOLODataset_binarized"
AUGMENTED_DATASET = "dataset/YOLODataset_augmented"

# === Augmentation Pipeline ===
augment = A.Compose([
    A.HorizontalFlip(p=0.5),
    A.VerticalFlip(p=0.3),
    A.RandomBrightnessContrast(p=0.5),
    A.Rotate(limit=30, p=0.5),
])

def make_dirs(base_dir):
    """Create images/train, images/val, labels/train, labels/val folders."""
    for subset in ["train", "val"]:
        os.makedirs(os.path.join(base_dir, "images", subset), exist_ok=True)
        os.makedirs(os.path.join(base_dir, "labels", subset), exist_ok=True)

# === Copy YAML file to new datasets ===
os.makedirs(BINARIZED_DATASET, exist_ok=True)
os.makedirs(AUGMENTED_DATASET, exist_ok=True)
shutil.copy(os.path.join(ORIGINAL_DATASET, "dataset.yaml"), BINARIZED_DATASET)
shutil.copy(os.path.join(ORIGINAL_DATASET, "dataset.yaml"), AUGMENTED_DATASET)

make_dirs(BINARIZED_DATASET)
make_dirs(AUGMENTED_DATASET)

def process_dataset(src_path, dst_bin_path, dst_aug_path):
    """Process images: binarize and augment, copy labels."""
    for subset in ["train", "val"]:
        src_img_dir = os.path.join(src_path, "images", subset)
        src_lbl_dir = os.path.join(src_path, "labels", subset)

        bin_img_dir = os.path.join(dst_bin_path, "images", subset)
        bin_lbl_dir = os.path.join(dst_bin_path, "labels", subset)
        aug_img_dir = os.path.join(dst_aug_path, "images", subset)
        aug_lbl_dir = os.path.join(dst_aug_path, "labels", subset)

        if not os.path.exists(src_img_dir):
            print(f"Folder not found: {src_img_dir}")
            continue

        for img_file in os.listdir(src_img_dir):
            img_path = os.path.join(src_img_dir, img_file)
            lbl_path = os.path.join(src_lbl_dir, os.path.splitext(img_file)[0] + ".txt")

            if not os.path.exists(lbl_path):
                print(f"Label not found for {img_file}, skipping.")
                continue

            # === Read original image ===
            image = cv2.imread(img_path)
            if image is None:
                print(f"Could not read image: {img_path}, skipping.")
                continue

            # === 1. Binarize ===
            gray = cv2.cvtColor(image, cv2.COLOR_BGR2GRAY)
            _, bin_image = cv2.threshold(gray, 127, 255, cv2.THRESH_BINARY)
            bin_image = cv2.cvtColor(bin_image, cv2.COLOR_GRAY2BGR)
            cv2.imwrite(os.path.join(bin_img_dir, img_file), bin_image)
            shutil.copy(lbl_path, os.path.join(bin_lbl_dir, os.path.basename(lbl_path)))

            # === 2. Augment ===
            transformed = augment(image=image)
            aug_image = transformed['image']
            cv2.imwrite(os.path.join(aug_img_dir, img_file), aug_image)
            shutil.copy(lbl_path, os.path.join(aug_lbl_dir, os.path.basename(lbl_path)))

if __name__ == "__main__":
    process_dataset(ORIGINAL_DATASET, BINARIZED_DATASET, AUGMENTED_DATASET)
    print("Datasets created: binarized and augmented.")

\end{lstlisting}
